% Created 2026-07-10 Fri 01:29
% Intended LaTeX compiler: pdflatex
\documentclass{article}
\usepackage[utf8]{inputenc}
\usepackage[T1]{fontenc}
\usepackage{graphicx}
\usepackage{longtable}
\usepackage{wrapfig}
\usepackage{rotating}
\usepackage[normalem]{ulem}
\usepackage{amsmath}
\usepackage{amssymb}
\usepackage{capt-of}
\usepackage{hyperref}

\usepackage{fancyhdr}
\usepackage[top=25mm, left=25mm, right=25mm]{geometry}
\usepackage{longtable}
\fancyhead[R]{}
\setlength\headheight{43.0pt}
\usepackage[T1]{fontenc}
\usepackage[utf8]{inputenc}
\usepackage[spanish, ]{babel}
\usepackage[backend=biber]{biblatex}

\author{Joset Muela, Dylan Zuñiga, Anthony Erazo}
\date{2026-07-10}
\title{Booteo de Sistema Operativo Linux}
\hypersetup{
 pdfauthor={Joset Muela, Dylan Zuñiga, Anthony Erazo},
 pdftitle={Booteo de Sistema Operativo Linux},
 pdfkeywords={},
 pdfsubject={},
 pdfcreator={Emacs 27.1 (Org mode 9.7.5)}, 
 pdflang={Español}}
\begin{document}

\maketitle
\fancyhead[C]{\includegraphics[scale=0.05]{../images/logoEPN.jpg}\\
ESCUELA POLITÉCNICA NACIONAL\\FACULTAD DE INGENIERÍA DE SISTEMAS\\
ARQUITECTURA DE COMPUTADORES}
\thispagestyle{fancy}




\section{Objetivos}
\label{sec:org6cb291e}

\begin{itemize}
\item Aprender a  bootear un sistema operativo Linux

\item Aprender sobre la secuencia de arranque de un computador como la
BIOS y la UEFI
\end{itemize}

\section{Instrucciones}
\label{sec:orgf4b4518}
\begin{enumerate}
\item Siga el tutorial para la creación de un booteable de Ubuntu como se
indica en la unidad 9 en el archivo Arranque Ubuntu\footnote{Dependiendo de las características de su computador debe
verificar si al crear el medio booteable para Ubuntu, su sistema
reconoce o acepta para el \emph{Target System}}.

\item Realice un respaldo del código de \emph{BitLocker} utilizando un command
prompt de Windows con permisos de administrador\footnote{Si no respalda el código, no podrá arrancar el disco luego de
reiniciar. Más información en: \href{https://www.partitionwizard.com/disk-recovery/bitlocker-recovery-key-bypass.html }{https://www.partitionwizard.com/disk-recovery/bitlocker-recovery-key-bypass.html}}:

\begin{verbatim}
    manage-bde-protectors C:-get
\end{verbatim}

\item Utilizando su computador o el del laboratorio, siga los pasos
indicados para obtener el booteo desde el USB con el Sistema
Operativo de Ubuntu

\item Realice el boot de Ubuntu desde la memoria USB externa y
seleccione la opción de \emph{Probar Sistema Operativo}.
\item Una vez ingresado active un terminal de Ubuntu: \texttt{Ctrl+ALT+T} y
obtenga las características del computador mediante la ejecución de
los comandos: \texttt{lscpu}, \texttt{lsb\_release -a} y \texttt{lspci}. Obtenga capturas
de la ejecución de los comandos. Luego abra este archivo en WSL y
ejecute los comandos desde los bloques shell. Para insertar un
bloque puede usar la combinación \texttt{C-c C-,}, escoja la opción \texttt{src}
e indique que es del tipo shell.
\end{enumerate}

\section{Desarrollo Emacs}
\label{sec:orgae4f420}
Se sugiere utilizar la opción \texttt{:fontsize
scriptsize} para dar formato a la ejecución del comando. Por ejemplo,
considere:

\begin{verbatim}
whoami
\end{verbatim}

\begin{verbatim}
leningfe
\end{verbatim}

\subsection{Comando \texttt{lsb\_release}}
\label{sec:org36af823}
\begin{verbatim}
    lsb_release -a
\end{verbatim}

\begin{center}
\begin{tabular}{ll}
Distributor ID: & Ubuntu\\[0pt]
Description: & Ubuntu 22.04.4 LTS\\[0pt]
Release: & 22.04\\[0pt]
Codename: & jammy\\[0pt]
\end{tabular}
\end{center}


\subsection{Comando \texttt{lspci}}
\label{sec:org9192909}
\begin{verbatim}
   lspci
\end{verbatim}

\begin{verbatim}
64cc:00:00.0 SCSI storage controller: Red Hat, Inc. Virtio 1.0 console (rev 01)
a731:00:00.0 System peripheral: Red Hat, Inc. Virtio file system (rev 01)
e689:00:00.0 3D controller: Microsoft Corporation Basic Render Driver
\end{verbatim}


\subsection{Comando \texttt{lscpu}}
\label{sec:org7ee376f}
\begin{verbatim}
   lscpu
\end{verbatim}

\begin{verbatim}
  Architecture:                            x86_64
  CPU op-mode(s):                          32-bit, 64-bit
  Address sizes:                           48 bits physical, 48 bits virtual
  Byte Order:                              Little Endian
  CPU(s):                                  2
  On-line CPU(s) list:                     0,1
  Vendor ID:                               AuthenticAMD
  Model name:                              AMD A9-9420 RADEON R5, 5 COMPUTE CORES 2C+3G
  CPU family:                              21
  Model:                                   112
  Thread(s) per core:                      1
  Core(s) per socket:                      2
  Socket(s):                               1
  Stepping:                                0
  BogoMIPS:                                5988.75
  Flags:                                   fpu vme de pse tsc msr pae mce cx8 apic sep mtrr pge mca cmov pat pse36 clflush mmx fxsr sse sse2 ht syscall nx mmxext fxsr_opt pdpe1gb rdtscp lm rep_good nopl cpuid extd_apicid tsc_known_freq pni pclmulqdq ssse3 fma cx16 sse4_1 sse4_2 movbe popcnt aes xsave avx f16c rdrand hypervisor lahf_lm cmp_legacy cr8_legacy abm sse4a misalignsse 3dnowprefetch osvw xop fma4 topoext ssbd ibpb vmmcall fsgsbase bmi1 avx2 smep bmi2 xsaveopt virt_ssbd arat
  Hypervisor vendor:                       Microsoft
  Virtualization type:                     full
  L1d cache:                               64 KiB (2 instances)
  L1i cache:                               192 KiB (2 instances)
  L2 cache:                                2 MiB (2 instances)
  NUMA node(s):                            1
  NUMA node0 CPU(s):                       0,1
  Vulnerability Gather data sampling:      Not affected
  Vulnerability Ghostwrite:                Not affected
  Vulnerability Indirect target selection: Not affected
  Vulnerability Itlb multihit:             Not affected
  Vulnerability L1tf:                      Not affected
  Vulnerability Mds:                       Not affected
  Vulnerability Meltdown:                  Not affected
  Vulnerability Mmio stale data:           Not affected
  Vulnerability Old microcode:             Not affected
  Vulnerability Reg file data sampling:    Not affected
  Vulnerability Retbleed:                  Mitigation; untrained return thunk; SMT disabled
  Vulnerability Spec rstack overflow:      Not affected
  Vulnerability Spec store bypass:         Mitigation; Speculative Store Bypass disabled via prctl
  Vulnerability Spectre v1:                Mitigation; usercopy/swapgs barriers and __user pointer sanitization
  Vulnerability Spectre v2:                Mitigation; Retpolines; IBPB conditional; STIBP disabled; RSB filling; PBRSB-eIBRS Not affected; BHI Not affected
  Vulnerability Srbds:                     Not affected
  Vulnerability Tsa:                       Not affected
  Vulnerability Tsx async abort:           Not affected
  Vulnerability Vmscape:                   Not affected
\end{verbatim}

\section{Desarrollo Capturas}
\label{sec:orga6a519e}
Coloque las capturas de los resultados obtenidos en los
terminales. 

\begin{figure}[htbp]
\centering
\includegraphics[width=.9\linewidth]{../images/lscpu1.jpeg}
\caption{\label{fig:orgd4d277c}Captura del comando lscpu - parte 1}
\end{figure}

\begin{figure}[htbp]
\centering
\includegraphics[width=.9\linewidth]{../images/lscpu2.jpeg}
\caption{\label{fig:org4d57647}Captura del comando lscpu - parte 2}
\end{figure}

\begin{figure}[htbp]
\centering
\includegraphics[width=.9\linewidth]{../images/lsb_realese.jpeg}
\caption{\label{fig:orgee7764f}Captura del comando lsb\textsubscript{release}}
\end{figure}

\begin{figure}[htbp]
\centering
\includegraphics[width=.9\linewidth]{..//images/lspci.jpeg}
\caption{\label{fig:orgf244c70}Captura del comando lspci}
\end{figure}


\clearpage

\section{Consulta}
\label{sec:org85d4d81}
Describa que hacen y para qué sirven los comandos utilizados en la práctica

\begin{description}
\item[{lscpu}] Muestra la información detallada sobre la arquitectura del procesador. Sirve para identificar características claves del hardware, fabricante, modelo, cantidad de núcleos y memoria caché, es muy útil para diagnosticar las capacidades del sistema.
\item[{lsb$\backslash$\textsubscript{release} -a}] Muestra información específica sobre la distribución de Linux que está instalada. El parámetro '-a' indica que imprima todos los detalles posibles.
\item[{lspci}] Es una herramienta que lista todos los buses PCI del computador y los dispositivos de hardware que están conectados a ellos. Sirve para identificar componentes físicos internos de la máquina, es fundamental para tareas de diagnóstico de hardware y búsqueda de controladores o drivers.
\end{description}
\end{document}
